\documentclass[runningheads]{llncs}

\usepackage[T1]{fontenc}
\usepackage{newtxtext}
\usepackage[varvw]{newtxmath}

\begin{document}

\title{Explaining Faults by Counterfactual Reasoning}
\author{Philipp Dirscherl}
\authorrunning{P. Dirscherl}
\institute{OTH Regensburg\\\email{your.email@example.com}}

\maketitle

\begin{abstract}
% Write this last. Summarize the question, the structural causal model,
% the published DC motor case, and the paper's theoretical conclusion.
\keywords{Counterfactual reasoning \and Fault responsibility \and
Model-based diagnosis}
\end{abstract}

\section{Introduction}
% Motivate symptom-level explanation when several faults are present.
% State the research question and delimit the paper as a theoretical analysis
% of published work. Use Reiter for the diagnosis background.

\section{Structural Models and Actual Causation}
% Explain the role of a causal model, context, and intervention.
% Present only the parts of Halpern and Pearl's actual-cause definition
% needed to understand the case: actuality, counterfactual dependence
% under a suitable contingency, and minimality.

\section{Counterfactual Fault Removal in a Dynamic System}
% Explain the time-indexed motor model at a conceptual level.
% State what is fixed between factual and counterfactual trajectories
% and what the closed-loop controller recomputes.
% Define symptom severity and the overheating event separately.

\section{Analysis of the Published DC Motor Case}
% Focus on Muji\'{c} and Wotawa's FEF4 + FMF1 experiment.
% Reproduce only the values from their Table 2 needed for the argument,
% with a citation in the caption or surrounding text.
% Explain why both faults affect severity but only FMF1 is necessary
% for their 90-degree-C overheating event in this simulated case.

\section{Discussion and Limitations}
% Develop the hypothetical 100-degree-C threshold as your own
% analytical deduction from the published values, not a new experiment.
% Explain dependence on event definition, model accuracy, and operating
% conditions. Discuss the authors' parameter-mismatch repetitions
% and their use of true simulated initial states.
% If using the formal term degree of responsibility, distinguish
% Chockler and Halpern's definition from the motor paper's score.

\section{Conclusion}
% Answer the research question for this case and state the scope
% of the conclusion in two or three sentences.

\begin{thebibliography}{4}

\bibitem{reiter1987}
Reiter, R.: A theory of diagnosis from first principles.
Artificial Intelligence \textbf{32}(1), 57--95 (1987).
\doi{10.1016/0004-3702(87)90062-2}

\bibitem{chockler2004}
Chockler, H., Halpern, J.Y.: Responsibility and blame:
a structural-model approach.
Journal of Artificial Intelligence Research \textbf{22}, 93--115 (2004).
\doi{10.1613/JAIR.1391}

\bibitem{halpern2005}
Halpern, J.Y., Pearl, J.: Causes and explanations:
a structural-model approach. Part I: Causes.
The British Journal for the Philosophy of Science
\textbf{56}(4), 843--887 (2005).
\doi{10.1093/bjps/axi147}

\bibitem{mujic2026}
Muji\'{c}, E., Wotawa, F.: Counterfactual fault responsibility
in a dynamic structural causal model of a DC motor.
In: 37th International Conference on Principles of Diagnosis
and Resilient Systems (DX 2026).
OASIcs, vol. 148, pp. 16:1--16:14.
Schloss Dagstuhl (2026).
\doi{10.4230/OASIcs.DX.2026.16}

\end{thebibliography}

\end{document}
